% ============================================================
% DRAFT NOTE (results)
% Contains: The one-shot confirmation result, falsification battery,
% stratified breakdowns, the quality-failure validation, and transfer.
% Written now: The frozen current-state-damage-v1 confirmation numbers
% (2026-09-01): per-horizon MSE vs action-clock, superiority bounds, and the
% falsification outcomes. Verify each number against
% results/recovered/current-state-damage-v1/confirmation_result/report.json
% before submission.
% Pending: Stratified tables, the quality-failure validation, and the
% transfer experiment (not yet run).
% Sources: current-state-damage-v1 confirmation result
% ============================================================
\section{Results}

\todo{STALE SECTION: written for the divergence objective; to be rewritten for the quality-risk target. Do not review yet.}
\label{sec:results}

\subsection{Forecasting near-term damage}

The confirmation set of 256 held-out prompt groups was opened and scored
exactly once, after every modeling decision was frozen on the 508 development
prompt groups. Table~\ref{tab:main-results} reports held-out MSE against the
action-clock comparator.

\begin{table}[t]
\centering
\caption{Held-out confirmation MSE by horizon. Improvement is relative MSE
reduction over the action-clock comparator. Simultaneous one-sided
superiority bounds on the paired loss difference are all below zero.}
\label{tab:main-results}
\begin{tabular}{lcccc}
\toprule
Horizon & HERALD & Action-clock & Improvement & Superiority bound \\
\midrule
$k{=}5$  & 0.006226 & 0.008386 & 25.8\% & $-0.001838$ \\
$k{=}10$ & 0.004349 & 0.006172 & 29.5\% & $-0.001561$ \\
$k{=}25$ & 0.002743 & 0.004244 & 35.4\% & $-0.001274$ \\
$k{=}50$ & 0.001988 & 0.003529 & 43.7\% & $-0.001225$ \\
\midrule
Macro    & 0.003826 & 0.005583 & 31.5\% & --- \\
\bottomrule
\end{tabular}
\end{table}

The causal signal strengthens with horizon: relative improvement grows from
25.8\% at five tokens to 43.7\% at fifty. \todo{Interpret: is this because
long-horizon damage is dominated by regime-level drift the causal features
detect early, or because short-horizon damage has irreducible per-token
noise? Check the development diagnostics.}

\subsection{Falsification battery}

The confirmation passed every preregistered check. Every task, compressor,
and ratio subgroup had a negative paired loss difference against the
comparator at every horizon; leave-one-task-out stress passed for all four
held-out tasks and horizons; simultaneous residual calibration intervals
contained zero and calibration-slope lower bounds were positive; predictions
were finite, nonnegative, and horizon-monotone; and an independent
recomputation reproduced the reported losses and superiority bounds.
\todo{Render the subgroup grid as a compact table or heatmap from the
confirmation report.}

\subsection{Where the signal holds and where it degrades}

\begin{table}[t]
\centering
\caption{Stratified held-out breakdown. \todo{Select the strata worth a
main-body table from the confirmation report; the full grid goes to the
appendix.}}
\label{tab:strata-results}
\begin{tabular}{llcccc}
\toprule
Stratum & Level & $k{=}5$ & $k{=}10$ & $k{=}25$ & $k{=}50$ \\
\midrule
\todo{task/compressor/ratio/position} & \todo{level} & \todo{v} & \todo{v} & \todo{v} & \todo{v} \\
\bottomrule
\end{tabular}
\end{table}

\subsection{Predicted divergence and realized quality failures}

% DRAFT NOTE: this experiment has NOT been run. It is the paper's
% load-bearing bridge from distribution-space damage to user-relevant
% failure; without it the motivation overreaches the evidence.
\todo{Run and report the preregistered secondary validation: per-run
aggregates of predicted damage versus realized task-quality failures and
catastrophe indicators.}
\begin{figure}[t]
\centering
\figplaceholder{separation of failed versus clean runs by per-run aggregate
of predicted damage, per task and compressor}{\todo{analysis output path}}
\caption{\todo{Quality-failure validation caption and interpretation.}}
\label{fig:quality-validation}
\end{figure}

\subsection{Cross-model transfer}

% DRAFT NOTE: not yet established; the frozen result explicitly excludes it.
\todo{Run and report zero-shot and adapted transfer on the second-family
traces once collected.}
\begin{table}[t]
\centering
\caption{Transfer to the second model family: zero-shot and after
lightweight adaptation.}
\label{tab:transfer-results}
\begin{tabular}{lcccc}
\toprule
Mode & $k{=}5$ & $k{=}10$ & $k{=}25$ & $k{=}50$ \\
\midrule
\todo{zero-shot} & \todo{v} & \todo{v} & \todo{v} & \todo{v} \\
\todo{adapted}   & \todo{v} & \todo{v} & \todo{v} & \todo{v} \\
\bottomrule
\end{tabular}
\end{table}
